\documentclass{article}

% if you need to pass options to natbib, use, e.g.:
%     \PassOptionsToPackage{numbers, compress}{natbib}
% before loading neurips_2024


% ready for submission
\usepackage[preprint]{neurips_2024}


% to compile a preprint version, e.g., for submission to arXiv, add add the
% [preprint] option:
%     \usepackage[preprint]{neurips_2024}


% to compile a camera-ready version, add the [final] option, e.g.:
%     \usepackage[final]{neurips_2024}


% to avoid loading the natbib package, add option nonatbib:
%    \usepackage[nonatbib]{neurips_2024}


\usepackage[utf8]{inputenc} % allow utf-8 input
\usepackage[T1]{fontenc}    % use 8-bit T1 fonts
\usepackage{hyperref}       % hyperlinks
\usepackage{url}            % simple URL typesetting
\usepackage{booktabs}       % professional-quality tables
\usepackage{amsfonts}       % blackboard math symbols
\usepackage{nicefrac}       % compact symbols for 1/2, etc.
\usepackage{microtype}      % microtypography
\usepackage{xcolor}         % colors


\title{Housing Price Prediction: A Comparative Study of Regularization Techniques}


% The \author macro works with any number of authors. There are two commands
% used to separate the names and addresses of multiple authors: \And and \AND.
%
% Using \And between authors leaves it to LaTeX to determine where to break the
% lines. Using \AND forces a line break at that point. So, if LaTeX puts 3 of 4
% authors names on the first line, and the last on the second line, try using
% \AND instead of \And before the third author name.


\author{%
  Miguel Nino Adalla (mva54)
  \And
  Reuben Thomas (rmt135)
  \And
  Nicole Krivokapic (nk752)
}


\begin{document}


\maketitle

\section{Executive Summary}

This project investigates linear regression models, enhanced through regularization, to predict housing prices using the Ames, Iowa Housing dataset. We systematically compared standard Linear Regression, Ridge Regression (L2), and Lasso Regression (L1), each trained on two datasets: one with and one without extensive feature engineering. The goal was to evaluate the impact of regularization and feature engineering on prediction accuracy.

The key findings include:
\begin{itemize}
    \item Ridge Regression with feature engineering was the best model, achieving R² = 0.937 and RMSE = 0.108.
    \item Feature engineering reduced test RMSE by up to 10\% and increased R² by up to 1.8\%.
    \item The most predictive features included ``TotalSF\_Log'', ``OverallQual'', ``GrLivArea\_Log'', and ``Neighborhood\_Price''.
    \item Regularization effectively addressed multicollinearity and improved generalization.
\end{itemize}

The developed models can assist various stakeholders in the real estate market, including homebuyers, sellers, and investors, by providing data-driven estimates of property values. Future work could explore more complex modeling approaches and incorporate additional external data sources to further improve prediction accuracy.

\section{Introduction}

This project addresses the problem of predicting housing prices in Ames, Iowa, through supervised learning. Specifically, we aim to determine how well linear models with regularization can estimate sale prices, and how feature engineering would influence their performance.

% To solve this problem, we implemented:
% \begin{enumerate}
%     \item Standard linear regression (Ordinary Least Squares)
%     \item Ridge regression (L2 regularization)
%     \item Lasso regression (L1 regularization)
% \end{enumerate}

% Each model was evaluated on both raw and feature-engineered datasets, for a total of six experiments. This approach allowed us to isolate the effect of feature engineering and regularization.

\section{Motivation}
This project is important because accurate housing price prediction helps various stakeholders make informed decisions. Potential homebuyers can determine if a house is fairly priced, sellers can set competitive prices, and investors can identify properties with growth potential. Understanding the factors that influence housing prices can inform urban planning and development strategies.

I'm particularly excited about this project because it provides an opportunity to apply statistical learning techniques to a real-world problem with tangible implications. Real estate is one of the largest markets in the economy, and data-driven approaches can help make this market more efficient and transparent.

Existing questions in this area include:
\begin{itemize}
    \item Which features have the strongest correlation with housing prices?
    \item How do different neighborhoods impact property values?
    \item Can regularization techniques improve prediction accuracy and model interpretability?
    \item Does feature engineering improve accuracy and interpretability?
\end{itemize}

\section{Method}

\subsection{Dataset}

For this project, I used the Ames Iowa Housing dataset, which contains information on residential properties sold in Ames, Iowa between 2006 and 2010. This comprehensive dataset includes 79 explanatory variables describing various aspects of residential homes. The data is in tabular format with both numerical and categorical features. Key feature categories include:
\begin{enumerate}
    \item Property characteristics: Square footage, number of rooms, bathrooms, etc.
    \item Location information: Neighborhood, zoning classification
    \item Quality and condition ratings: Overall quality, exterior condition, etc.
    \item Time-related features: Year built, year remodeled, year sold
    \item Special features: Fireplaces, pools, garages, etc.
\end{enumerate}

% \subsection{Problem Definition and Feature Space}

% The problem is defined as a regression task where we predict a continuous value (SalePrice) based on property features. The feature space was refined through:
% \begin{enumerate}
%     \item Data preprocessing: Handling missing values, outlier detection, and transformations
%     \item Feature engineering: Creating new features, encoding categorical variables
% \end{enumerate}

\subsection{Problem Definition \& Model Selection}

Because during EDA we observed that SalePrice was slightly skewed right, we chose to add a new log-transformed feature, SalePrice\_Log, to use as the target variable. Our goal was to train three regression models to predict SalePrice\_Log. We implemented these three models, each on two datasets:
\begin{enumerate}
    \item Standard Linear Regression: Provides a baseline performance using Ordinary Least Squares (OLS)
    \item Ridge Regression: Applies L2 regularization to reduce model complexity and mitigate multicollinearity
    \item Lasso Regression: Employs L1 regularization which performs feature selection by shrinking less important feature coefficients to zero
\end{enumerate}

Each model was evaluated on both raw and feature-engineered datasets, for a total of six experiments. This approach allowed us to isolate the effect of feature engineering and regularization.

\subsection{Implementation Steps}

\begin{enumerate}
    \item Exploratory Data Analysis (EDA):
    \begin{itemize}
        \item Analyze feature distributions
        \item Identify correlations between features and target variable
        \item Detect any outliers and anomalies, if any
    \end{itemize}
    \item Data Preprocessing:
    \begin{itemize}
        \item Handle missing values through imputation or removal
        \item Transform skewed numerical features using log or other transformations
        \item Encode categorical variables using one-hot encoding or label encoding
    \end{itemize}
    \item Feature Engineering:
    \begin{itemize}
        \item Create new features that might have predictive power
        \item Apply feature scaling to normalize the range of features
    \end{itemize}
    \item Model Training:
    \begin{itemize}
        \item Split data into training and testing sets
        \item Train three different linear regression models each with both dataset types
        \item Tune hyperparameters using cross-validation
    \end{itemize}
    \item Model Evaluation:
    \begin{itemize}
        \item Assess performance on test set
        \item Compare models using various metrics
        \item Analyze the improvement in using feature engineering for each of the three models
    \end{itemize}
    \item Predictions Visualizations:
    \begin{itemize}
        \item Analyze predictions of the best performing model
        \item Visualize relationships between actual and predicted sales prices
        \item Visualize residual distribution between actual and predicted sales prices
    \end{itemize}
    \item Feature Importance:
    \begin{itemize}
        \item Compare feature importance across different models
        \item Visualize the most influential features using bar charts
    \end{itemize}
    \item Conclusion
\end{enumerate}

% \subsection{Feature Engineering}

% Feature engineering was a critical component of our housing price prediction methodology, significantly enhancing model performance. The created features include:

% \subsubsection{Time-Based Features}
% \begin{itemize}
%     \item HouseAge, YearsSinceRemodel, IsRemodeled
% \end{itemize}

% \subsubsection{Area and Size Features}
% \begin{itemize}
%     \item TotalSF: total square footage
%     \item TotalBath: Weighted sum of bathrooms
%     \item TotalPorchSF: Combined square footage of all porch types
% \end{itemize}

% \subsubsection{Ratio and Proportion Features}
% \begin{itemize}
%     \item LotRatio: Calculated as LotFrontage / LotArea
%     \item BedroomRatio: Ratio of bedrooms to total rooms
% \end{itemize}

% \subsubsection{Interaction and Indicators}
% \begin{itemize}
%     \item OverallQualCond: Multiplication of OverallQual and OverallCond
%     \item HasPool, HasGarage, HasFireplace
% \end{itemize}

% \subsubsection{Neighborhood Encoding}
% \begin{itemize}
%     \item Neighborhood\_Price: Mean sale price by neighborhood
% \end{itemize}

\subsection{Evaluation Metrics}

To evaluate model performance, we used the root mean squared error (RMSE), mean absolute error (MAE), R squared score (R2), cross-validation scores (5 fold).
% \begin{itemize}
%     \item Root Mean Squared Error (RMSE)
%     \item Mean Absolute Error (MAE)
%     \item R-squared score (R²)
%     \item Cross-validation scores (5-fold)
% \end{itemize}

\section{Insights and Methodology Rationale}

\subsection{Exploratory Data Analysis Insights}

\subsubsection{Target Variable Transformation}
We chose to log-transform the SalePrice target variable because the original distribution exhibited significant right skewness (1.74 skewness value). This transformation helped reduce this positive skewness, and create a more normal distribution. This is important for linear models which assume that the target variable is normally distributed.

\subsubsection{Key Correlation Findings}
The correlation analysis revealed several features with strong positive correlations to SalePrice, namely OverallQual, GrLivArea, GarageCars, and TotalBsmtSF. This helped us determine which features will likely be most predictive and influential for the models. It also helped guide our feature engineering efforts, particularly in creating composite features that amplified these relationships.
% \begin{itemize}
%     \item OverallQual: Demonstrated the highest correlation, indicating the subjective quality assessment is extremely predictive
%     \item GrLivArea: Above-ground living area showed strong correlation, confirming the importance of home size
%     \item GarageCars and TotalBsmtSF: Both showed substantial correlation, highlighting the value placed on these amenities
% \end{itemize}
% These findings guided our feature engineering efforts, particularly in creating composite features that amplified these relationships.

\subsubsection{Categorical Feature Analysis}
The boxplot analysis of categorical variables like Neighborhood and ExterQual revealed significant group-level pricing patterns. Some neighborhoods consistently showed higher median prices, while quality metrics for kitchens and exteriors demonstrated substantial price impact across quality levels. This confirmed the need for effective encoding strategies for these categorical variables.

\subsection{Preprocessing Rationale}

\subsubsection{Missing Value Strategy}
We needed to rigorously handle missing values because our models would not be able to handle them if kept in the dataset. At the same time, we didn't want to sacrifice any critical information. Therefore, our missing value imputation approach was carefully tailored to each feature type.
\begin{itemize}
    \item For categorical features like FireplaceQu, missing values were replaced with 'None' to preserve the semantic meaning (absence of a fireplace)
    \item For numerical features like LotFrontage, we employed group-based imputation using neighborhood medians, providing more accurate estimates than global measures
\end{itemize}
This approach prevented information loss while maintaining data integrity for model training

\subsubsection{Feature Transformation \& Encoding}
We applied log transformations to heavily right-skewed features like LotArea, MiscVal, and BsmtFinSF1. By doing so we make their relationships more linear, reduce the influence of outliers, and ultimately improve the performance of our linear models.
Our encoding approach was feature-specific:
\begin{itemize}
    \item Ordinal encoding was applied to quality features (ExterQual, KitchenQual, etc.) to capture their inherent ordering (Excellent > Good > Average, etc.)
    \item One-hot encoding was used for nominal categorical features with low cardinality, allowing models to learn group differences without imposing artificial ordering
    \item These encoding choices allowed our models to effectively utilize both ordered and unordered categorical information
\end{itemize}

\subsection{Feature Engineering Rationale}

\subsubsection{Domain-Informed Transformations}
Our feature engineering decisions were driven by domain knowledge of real estate valuation and by our prior observations obtained during EDA. Some important new features include:
\begin{itemize}
    \item Time-Based Features (eg. HouseAge, IsRemodeled): Captures potential value added through renovation or lost through depreciation. These are often not linearly related to age.
    \item Area and Size Features (eg. TotalSF, TotalBath): Created to provide a more comprehensive size metric than any individual area measurement.
    \item Ratio and Proportion Features (e.g. LotRatio, BedroomRatio): Captures relationship between more relative relationships that might be missed by linear models. 
    \item Products and Aggregation of Features (eg. OverallQualCond, Neighborhood\_Price): Captures more general aspects of homes that could have a combined effect beyond their individual impact.
\end{itemize}

\subsection{Model Selection and Hyperparameter Tuning}

\subsubsection{Regularization Approach}
We specifically chose Ridge and Lasso regression models to address challenges in the housing dataset:
\begin{itemize}
    \item The high-dimensional feature space (310 features after engineering) created overfitting risks
    \item Multicollinearity between features (e.g., various area measurements) could destabilize standard linear regression
    \item Ridge (L2) regularization shrinks all coefficients to reduce model complexity
    \item Lasso (L1) regularization performs feature selection by reducing some coefficients to zero
\end{itemize}

\subsubsection{Hyperparameter Selection}
Our GridSearch explored a range of alpha values:
\begin{itemize}
    \item Ridge: [0.001, ..., 1000]
    \item Lasso: [0.0001, ..., 10]
\end{itemize}
This wide range was deliberately chosen to explore the spectrum from under-regularization to over-regularization, allowing us to identify the optimal balance between model fit and generalization.

\section{Results}

\subsection{Model Performance}

After training and evaluating all the models, I found that the regularized models performed better than the standard linear regression. Ridge regression with feature engineering showed the best overall performance with an RMSE of 0.108 and R² of 0.937, closely followed by Lasso regression (FE) with similar metrics.

\subsection{Performance Metrics Analysis}

\begin{table}[htbp]
\centering
\begin{tabular}{lcccc}
\toprule
Model & RMSE (Test) & MAE (Test) & R² (Test) & CV RMSE \\
\midrule
Linear Regression (No FE) & 0.1330 & 0.0905 & 0.9043 & 0.1582 \\
Ridge Regression (No FE) & 0.1182 & 0.0798 & 0.9245 & 0.1407 \\
Lasso Regression (No FE) & 0.1202 & 0.0807 & 0.922 & 0.1435 \\
Linear Regression (FE) & 0.1161 & 0.0830 & 0.9270 & 0.1418 \\
Ridge Regression (FE) & 0.1076 & 0.0749 & 0.9374 & 0.128 \\
Lasso Regression (FE) & 0.108 & 0.0755 & 0.937 & 0.1269 \\
\bottomrule
\end{tabular}
\caption{Performance metrics across all models}
\end{table}

Feature Engineering Benefits:
\begin{itemize}
    \item Ridge: $\sim$9\% RMSE reduction, +1.4\% R²
    \item Lasso: $\sim$10\% RMSE reduction, +1.6\% R²
    \item Linear: $\sim$13\% RMSE reduction, +2.5\% R²
\end{itemize}

% \subsection{Visualizations}

% The visualizations from the analysis reveal several key insights about housing price prediction in Ames, Iowa:
% \begin{enumerate}
%     \item The distribution of sale prices presented by histogram showed that Sales Price was positively skewed. This was addressed with log-transformation.
%     \item The multicollinearity heatmap and the scatter plots confirmed relationships between sale price and overall quality, garage area, and Total Basement sqft.
%     \item Feature importance plots highlight that quality/condition metrics (e.g. overall quality, kitchen quality, sale condition) and size-related variables (e.g. TotalSF, GrLivArea) contribute most significantly to price determination.
% \end{enumerate}

\subsection{Prediction Visualization Analysis}

\subsubsection{Actual vs. Predicted Analysis}
The comparison of actual vs. predicted values showed that the best model (Ridge Regression with feature engineering) demonstrated very tight clustering along the identity line, indicating very strong predictive accuracy. The model performed consistently across most of the price range. However there looked to be a slight tendency to under predict for the highest-value homes, suggesting potential for improvement in this region.

\subsubsection{Residual Analysis}
The residual plots provided additional diagnostics. The scatterplot confirmed that the model was very consistent, and that it had a tendency to underpredict high-value properties. The distribution was approximately normal with a very slight right skew, again confirming the underprediction tendency.
% \begin{itemize}
%     \item The residuals showed no strong patterns across the prediction range, supporting the assumption of homoscedasticity
%     \item The residual distribution was approximately normal with a slight tail, confirming some underprediction for high-value properties
%     \item The lack of distinct patterns in the residuals suggested the model captured most systematic variation in the data
% \end{itemize}

\subsection{Feature Importance Insights}

The feature importance analysis identified the most influential predictors:
\begin{itemize}
    \item TotalSF\_Log: Emerged as the single most important feature, confirming that total home size is the primary value driver
    \item OverallQual: The subjective quality rating showed very high importance, highlighting the value of expert assessment
    \item GrLivArea\_Log: Above-ground living area remained highly influential even with total square footage present
    \item Neighborhood\_Price: This engineered feature effectively captured location value, demonstrating the success of our feature engineering approach
\end{itemize}

    Between the the Ridge and Lasso models, both approaches identified similar top features. TotalSF\_Log was the most important feature in both models, showing the strongest positive coefficient. However, there were some interesting differences. Lasso tended to prefer a few, very important features, whereas Ridge regression distributed importance across more features. This is characteristic of L1 and L2 regularization, and confirms that our models behaved normally.

\section{Discussion}

\subsection{Expected vs. Actual Results}

We expected that both Ridge and Lasso regression would outperform OLS due to their ability to handle multicollinearity and prevent overfitting. We anticipated that Lasso might perform particularly well given its feature selection capability in a dataset with numerous features.

Our actual results largely aligned with these expectations. Both regularized models outperformed the standard linear regression, confirming the value of regularization for this problem. However, Ridge regression slightly outperformed Lasso regression. This suggests that feature shrinkage (reducing the magnitude of all coefficients) was more beneficial than feature selection (eliminating some features entirely) for this particular dataset.

The feature importance analysis revealed some surprising insights. While we expected features like overall quality and above-ground living area to be significant predictors, we did not anticipate the substantial impact of basement-related features. The finished basement area emerged as one of the top predictors, indicating that below-ground living space is highly valued in this market.

% Additionally, I was surprised by the relatively modest impact of lot size compared to home quality metrics. This suggests that in Ames, Iowa, homebuyers prioritize the quality and condition of the structure more than the size of the land it sits on. This finding contrasts with studies from other real estate markets where land value is often a dominant factor, particularly in high-density urban areas.

% \subsection{Impact of Feature Engineering}

% Our feature engineering efforts expanded the feature space from the original set to a significantly larger number of engineered features (multiple new features were created). This process:
% \begin{enumerate}
%     \item Captured domain knowledge: By incorporating real estate principles into features like age, total area, and quality-condition interactions.
%     \item Improved model performance: The R² score increased from 0.89 with basic features to 0.94 with engineered features.
%     \item Enhanced interpretability: Created features that align with intuitive understanding of housing valuation factors.
%     \item Addressed statistical requirements: Transformations helped meet linear regression assumptions by normalizing distributions and reducing multicollinearity.
% \end{enumerate}

\subsection{Limitations and Future Work}

Despite the strong performance of our models, several limitations should be acknowledged:
\begin{enumerate}
    \item Temporal relevance: The dataset only includes houses sold between 2006-2010, potentially limiting generalizability to current market conditions, especially considering the significant fluctuations in the housing market since that period.
    \item Geographic specificity: The models are trained on data from Ames, Iowa only and may not transfer well to other real estate markets with different valuation dynamics.
    \item Missing external factors: Economic indicators, interest rates, and market trends are not captured in the dataset but likely influence housing prices significantly.
    \item Model simplicity: While linear models offer interpretability, they may miss complex non-linear relationships between features and price.
\end{enumerate}

To improve this analysis in future work, I would recommend:
\begin{enumerate}
    \item Advanced modeling techniques: Implement ensemble methods or gradient boosting models to capture more complex patterns and potentially improve prediction accuracy.
    \item Time series analysis: Incorporate temporal trends in housing prices to account for market fluctuations and seasonality effects.
    \item External data integration: Include economic indicators, school district ratings, crime statistics, and neighborhood development data to provide additional context.
    \item Neural network approach: Explore deep learning techniques for potentially higher accuracy, especially for capturing non-linear relationships in the data.
\end{enumerate}

\section{Conclusion}

In conclusion, this project demonstrated the effectiveness of regularized linear regression techniques and feature engineering for housing price prediction. Both Ridge and Lasso provided improvements over standard linear regression, with Ridge regression showing the best overall performance. The analysis highlighted the importance of overall quality, living area size, basement features, and neighborhood in determining housing prices, providing valuable insights for real estate stakeholders, homebuyers, and sellers in the Ames market.

The regularization approach successfully addressed multicollinearity issues common in housing data and prevented overfitting. The interpretability of the linear models allowed for clear identification of key value drivers, making the findings actionable for various stakeholders in the real estate market.

Future research directions could explore more complex modeling techniques and incorporate additional data sources to further enhance prediction accuracy, particularly for luxury properties where the current models showed slightly lower performance.

\end{document}